\section{Related Work}
\label{sec:related}

The contribution of this paper sits at the intersection of four research
threads. Gradient-trained IoT intrusion detection establishes the
closed-set accuracy ceiling. Large language models applied to network
security expose the inference-cost gap that motivates an
offline-only architecture. Edge deployment of foundation models supplies
the SMCC vocabulary that frames our memory-computation split. Trustworthy
detection along the zero-day, adversarial, and privacy axes defines the
operating envelope the paper measures. We review each thread in turn and
close by naming the gap that no prior work jointly fills.

\subsection{Deep Learning and Tree-Based IDS for IoT}
\label{subsec:rw_dl_ids}

Gradient-trained classifiers define the accuracy bar that any
deployable IoT detector must clear. Random Forest and decision trees
remain the dominant baselines on every benchmark we evaluate because they
exploit gradient-optimised splits over the same numeric flow features
that the underlying data exposes. The five datasets used in this paper
are the canonical IoT benchmarks introduced by Neto~et~al.\
\cite{neto2023ciciot2023} for CIC-IoT2023, Alsaedi~et~al.\
\cite{alsaedi2020toniot} for TON\_IoT, Koroniotis~et~al.\
\cite{koroniotis2019botiot} for Bot-IoT, Moustafa and Slay
\cite{moustafa2015unswnb15} for UNSW-NB15, and Zolanvari~et~al.\
\cite{zolanvari2019wustl,zolanvari2021wustl} for WUSTL-IIoT. Recent deep
detectors push these benchmarks toward saturation. Transformer-based
multi-class classification on CIC-IoT-2023~\cite{wang2024transformer},
LSTM detectors on the same
benchmark~\cite{aboulela2025transformer,deeplearning2025lstm}, and
federated dataset-centric evaluations~\cite{federated2025iotids} all
report near-perfect macro-F1 on cleanly separable classes. The common
constraint is that each new dataset requires fresh gradient training,
fresh feature scaling, and a new model artefact that must coexist with
the gateway's routing workload. Our work targets the same accuracy bar
without per-dataset gradient training and without shipping model weights
to the gateway.

\subsection{Large Language Models and RAG for Network Intrusion Detection}
\label{subsec:rw_llm_ids}

The closest prior work uses an LLM directly inside the detection
pipeline. RuleMaster+ fine-tunes an instruction-tuned LLM to emit
Snort-style IDS rules and reports gains over zero-shot
ChatGPT~\cite{lian2024rulemaster}. Hybrid LLM frameworks combine a
small LLM with a numerical classifier to handle zero-day IoT
traffic~\cite{alhammouri2025hybridllm,sensors2026hybrid_llm_iiot}. Agentic
designs such as IDS-Agent reason over tool outputs at query
time~\cite{idsagent2024}. Lightweight LLMs for on-device attack
detection compress an LLM small enough to run on a single
gateway~\cite{comcomap2025lightweight}, and broader evaluations of LLMs
on IoT IDS benchmarks confirm that fine-tuned models match Random Forest
within a few F1
points~\cite{aboulela2025evaluating,iot2025llm_ecosystem_survey}.
Retrieval-augmented variants extend the same line. RLFE-IDS pairs a
fine-tuned LLM with a custom embedding model and a retrieval
store~\cite{li2025rlfeids}, and AgCyRAG combines knowledge-graph and
vector retrieval for security analysis~\cite{agcyrag2025}. Both inherit
the foundational RAG formulation of Lewis~et~al.\
\cite{lewis2020rag} and use Sentence-BERT-style dense embeddings
\cite{reimers2019sbert,song2020mpnet}.

Every approach above pays an LLM-inference cost at the moment a flow
arrives. Fine-tuned and distilled variants reduce the per-call latency
but retain parametric weights on the device. Retrieval-augmented
variants further add an embedding step and a vector lookup to the
per-query path. None of these designs eliminate the model artefact at
inference time, which is the constraint a single-digit-watt IoT gateway
cannot relax. Table~\ref{tab:rw_compare} summarises this distinction.
Our architecture is the only entry that combines an offline LLM stage
with a stateless rule evaluator at runtime, and the only one that
exposes retrieval confidence as an explicit zero-day score rather than
as an implicit byproduct of generation.

\begin{table*}[t]
  \centering
  \caption{Comparison of representative LLM-assisted IDS approaches.
    \emph{LLM stage} indicates whether the language model runs offline at
    policy synthesis time or per query at inference time.
    \emph{Runtime artefact} names the object deployed on the gateway.
    \emph{Edge-deployable} marks designs whose runtime cost fits inside a
    single-digit-watt IoT gateway envelope.
    \emph{Zero-day} marks designs that expose an explicit
    out-of-distribution score.}
  \label{tab:rw_compare}
  \begin{tabular}{lllll}
    \toprule
    \textbf{Approach} &
    \textbf{LLM stage} &
    \textbf{Runtime artefact} &
    \textbf{Edge-deployable} &
    \textbf{Zero-day} \\
    \midrule
    RuleMaster+~\cite{lian2024rulemaster}              & Offline (rule gen.)   & Snort rules + LLM updates    & Partial   & No          \\
    Hybrid LLM-NIDS~\cite{alhammouri2025hybridllm}          & Per-query             & Frozen LLM + classifier      & No        & Yes         \\
    Sensors hybrid IIoT~\cite{sensors2026hybrid_llm_iiot} & Per-query         & Frozen LLM encoder + MLP     & No        & No          \\
    IDS-Agent~\cite{idsagent2024}                     & Per-query (agentic)   & LLM agent + tools            & No        & Yes         \\
    RLFE-IDS~\cite{li2025rlfeids}                       & Per-query             & Fine-tuned LLM + RAG store   & No        & Score-based \\
    AgCyRAG~\cite{agcyrag2025}                        & Per-query             & LLM + KG + vector store      & No        & No          \\
    Lightweight LLM~\cite{comcomap2025lightweight}    & Per-query             & Compressed LLM               & Partial   & No          \\
    \textbf{Ours}                                     & \textbf{Offline only} & \textbf{JSON threshold rules} & \textbf{Yes} & \textbf{Score-based (SMCC)} \\
    \bottomrule
  \end{tabular}
\end{table*}

\subsection{Edge Deployment of Foundation Models and the SMCC Paradigm}
\label{subsec:rw_smcc}

Foundation models cannot run at line rate on an IoT gateway under any
known compression budget. Surveys of integrated sensing, communication,
and computation formalise this constraint and propose resource sharing
across the three functionalities as the path to deployable edge
intelligence~\cite{iscc2024survey,wen2023otaiscc,saadat2025iscc6g}.
Edge large AI model work pushes the same point in a foundation-model
register and identifies the offline-online split as the dominant lever
for practical deployment~\cite{wang2025edgelam,frontiers2025edgellm}.
Edge AI surveys treat model compression, pruning, and distillation as
necessary but insufficient~\cite{deploying2025edgeai}. The Integrated
Sensing, Memory, Communication, and Computation extension that frames
this special issue adds explicit memory as a first-class stage and
gives compressed offline reasoning a place in the architecture. Our
design instantiates this split. The retrieval index and the LLM act
together as the memory stage that is touched only at training time. The
deployed rule evaluator carries no model weights and runs as the
computation stage at line rate.

\subsection{Trustworthy IDS Across Zero-Day, Adversarial, and Privacy Axes}
\label{subsec:rw_trust}

Zero-day detection in network security has converged on two approaches.
Unsupervised anomaly scoring identifies traffic that deviates from a
learned normal distribution~\cite{guo2023zerodaysurvey,sarhan2023zeroshot},
and explainable models flag novel attacks through attention or feature
attribution~\cite{discoveriot2025zeroday_xai,kasongo2025attnzeroday}. Both
families assume access to a trained discriminative model that produces a
calibrated confidence. Our retrieval-confidence score is a third option.
It uses the neighbour distribution in the offline vector store as the
out-of-distribution signal, which makes the score available without
training a separate anomaly detector.

Adversarial robustness work on network IDS frames evasion under the
Kerckhoffs assumption that the attacker knows the
detector~\cite{apruzzese2021realistic,andrew2024adversarialchallenges,khamaiseh2023evasionreview}.
The canonical evaluation target is Kitsune~\cite{mirsky2018kitsune}, and
recent functionality-preserving attacks confirm that threshold-based
detectors fall to small budgets when a single feature governs the
boundary~\cite{teuffenbach2022funcpreserving}. Our Kerckhoffs sweep
extends this evaluation to LLM-generated rules and reports
$\varepsilon$@ESR$=0.5$ alongside trained tree baselines, showing that
the LLM policy distributes evasion cost across features rather than
concentrating it on a bottleneck.

Privacy-preserving IDS work attacks the orthogonal axis. Differential
privacy on network traces protects shared
data~\cite{fan2021condensationdp}, federated learning with feature
reduction protects participating
clients~\cite{springer2025federatedfeatured}, and homomorphic encryption
combined with graph neural networks protects inference
itself~\cite{scireports2025hegnn_iiot}. None of these methods address
the leakage channel that the prompt itself opens when feature names
encode topology or vendor information. Our anonymization ablation
measures that channel directly and shows that opaque identifiers do not
degrade detection.

Rule-based and explainable IDS designs supply the deployment archetype
this paper inherits. Recent rule-induction frameworks
\cite{sensors2025ruleinduction} and self-attention explainable IoT
detectors~\cite{scireports2025xaiids} confirm that interpretable
threshold rules remain competitive when training data is well separated.
Our work differs in that the rules are synthesised by an LLM from a
retrieved sample rather than induced from the full training set.

\subsection{Positioning}
\label{subsec:rw_position}

No prior work jointly delivers foundation-model-quality rule synthesis,
zero LLM cost at inference, an ACL-compatible runtime artefact, a
retrieval-confidence zero-day extension, a Kerckhoffs robustness audit,
and a feature-name anonymization study. The closest LLM-IDS designs pay
an inference-time LLM cost the gateway cannot absorb. The closest
zero-day designs assume a discriminative confidence we deliberately do
not train. The closest robustness studies evaluate trained classifiers
rather than LLM-emitted rules. The closest privacy studies target trace
release rather than prompt leakage. The SMCC memory-computation split
is the architectural lens that lets a single design address all four
axes at once, and the remainder of this paper realises that design.
